\documentclass[11pt]{article}
\usepackage[a4paper,margin=18mm]{geometry}
\usepackage[no-math]{fontspec}
\setmainfont{Latin Modern Roman}
\usepackage{amsmath,amssymb,amsthm,xfrac,xspace,hyperref,graphicx,comment}
\excludecomment{DefTralics}

\providecommand{\bibliofont}{\small}
\newcommand{\ie}{i.e.,\xspace}
\newcommand{\eg}{e.g.,\xspace}
\newcommand{\etc}{etc.\xspace}
\newcommand{\cf}{cf.\xspace}
\newcommand{\resp}{resp.\xspace}
\newcommand{\smaller}{\small}
\newcommand{\Subsection}{\subsection}
\newcommand{\Subsubsection}{\subsubsection}
\newcommand{\skpt}{\smallskip}
\emergencystretch=3em
\numberwithin{equation}{section}\input{author-preamble.tex}
\makeatletter
\renewcommand\section{\@startsection{section}{1}{\z@}{-3.5ex \@plus -1ex \@minus -.2ex}{2.3ex \@plus .2ex}{\normalfont\large\bfseries}}
\makeatother
\begin{document}
\begin{center}{\Large Stable item 2046}\end{center}
\paragraph{Source and attribution.}
Andriy Regeta, Immanuel van Santen, \emph{Characterizing smooth affine spherical varieties via the automorphism group}, Journal de l’École polytechnique — Mathématiques 8 (2021), publisher version, DOI 10.5802/jep.149. \href{https://jep.centre-mersenne.org/articles/10.5802/jep.149/}{Publisher article}. Authors' text: \href{https://creativecommons.org/licenses/by/4.0/}{CC BY 4.0}.
\paragraph{Editorial problem boundary.}
Prove only Proposition 7.3 below over an algebraically closed characteristic-zero field, for an irreducible quasi-affine variety with faithful action of a connected solvable group containing nontrivial unipotent elements. The open-orbit condition is equivalent to the printed bounded semi-invariant vector-field spaces and bounded generalised root subgroups. The complete equal-weight semi-invariant and arbitrarily large faithful replicated additive-action construction is retained. Reductive reconstruction, arbitrary torus actions and later weight-monoid formulae are context only.
\paragraph{Difficulty (editorial): H2.}
The cited rational quotient and solvable invariant-fraction results do not by themselves produce root subgroups of a prescribed common weight and arbitrarily large dimension. The author combines two equal-weight semi-invariants with a homogeneous additive action to construct the displayed replicated actions, and verifies faithfulness through the independence of their homogeneous powers. That same-weight replication construction supplies the unbounded family needed for the converse implication, together with the exact vector-field/root-subgroup comparison.
\paragraph{Editorial transcription note.}
Original author language, mathematics and ordering are preserved. Publisher front matter is replaced by this independent article wrapper. Bibliography is recovered from matching cached publisher HTML, including its TeX formulas. No new proof or mathematical repair is supplied. Surrounding original results are retained for conservative context and are not extra selected corpus claims.
\clearpage
\input{author-body.tex}
\begin{thebibliography}{99}
\bibitem{ArGa2010Cox-rings-semigrou} I. V. Arzhantsev \& S. A. Gaĭfullin - “Cox rings, semigroups, and automorphisms of affine varieties” , Mat. Sb. (N.S.) 201 (2010) no. 1, p. 3-24
\bibitem{AuTe2003Asymptotic-cones-a} A. Auslender \& M. Teboulle - Asymptotic cones and functions in optimization and variational inequalities , Springer Monographs in Math. , Springer-Verlag, New York, 2003
\bibitem{BrCu2010Classification-of-} P. Bravi \& S. Cupit-Foutou - “Classification of strict wonderful varieties” , Ann. Inst. Fourier (Grenoble) 60 (2010) no. 2, p. 641-681
\bibitem{Be2003Lifting-of-morphis} F. Berchtold - “Lifting of morphisms to quotient presentations” , Manuscripta Math. 110 (2003) no. 1, p. 33-44
\bibitem{BrPe2005Wonderful-varietie} P. Bravi \& G. Pezzini - “Wonderful varieties of type $D$” , Represent. Theory 9 (2005), p. 578-637
\bibitem{Br2007Wonderful-varietie} P. Bravi - “Wonderful varieties of type $E$” , Represent. Theory 11 (2007), p. 174-191
\bibitem{Br2010Introduction-to-ac} M. Brion - “Introduction to actions of algebraic groups” , in Actions hamiltoniennes: invariants et classification , vol. 1 , Centre Mersenne, Grenoble, 2010, p. 1-22, https://ccirm.centre-mersenne.org/volume/CCIRM\_2010\_\_1/
\bibitem{Cu2014Wonderful-Varietie} S. Cupit-Foutou - “Wonderful varieties: a geometrical realization” , 2014
\bibitem{CoLiSc2011Toric-varieties} D. A. Cox, J. B. Little \& H. K. Schenck - Toric varieties , Graduate Studies in Math. , vol. 124 , American Mathematical Society, Providence, RI, 2011
\bibitem{CaReXi2019Families-of-Commut} S. Cantat, A. Regeta \& J. Xie - “Families of commuting automorphisms, and a characterization of the affine space” , 2019
\bibitem{De1970Sous-groupes-algeb} M. Demazure - “Sous-groupes algébriques de rang maximum du groupe de Cremona” , Ann. Sci. École Norm. Sup. (4) 3 (1970), p. 507-588
\bibitem{FuKr2018On-the-geometry-of} J.-P. Furter \& H. Kraft - “On the geometry of the automorphism groups of affine varieties”
\bibitem{Fr2017Algebraic-theory-o} G. Freudenburg - Algebraic theory of locally nilpotent derivations , Encyclopaedia of Math. Sciences , vol. 136 , Springer-Verlag, Berlin, 2017
\bibitem{Fu1993Introduction-to-to} W. Fulton - Introduction to toric varieties , Annals of Math. Studies , vol. 131 , Princeton University Press, Princeton, NJ, 1993
\bibitem{Gr1961Elements-de-geomet-II} A. Grothendieck - “Éléments de géométrie algébrique. II. Étude globale élémentaire de quelques classes de morphismes” , Publ. Math. Inst. Hautes Études Sci. 8 (1961), p. 1-222
\bibitem{Gr1997Algebraic-homogene} F. D. Grosshans - Algebraic homogeneous spaces and invariant theory , Lect. Notes in Math. , vol. 1673 , Springer-Verlag, Berlin, 1997
\bibitem{Hu1975Linear-algebraic-g} J. E. Humphreys - Linear algebraic groups , Graduate Texts in Math. , vol. 21 , Springer-Verlag, New York-Heidelberg, 1975
\bibitem{Ka2005On-a-theorem-of-Ax} S. Kaliman - “On a theorem of Ax” , Proc. Amer. Math. Soc. 133 (2005) no. 4, p. 975-977
\bibitem{Kn1991The-Luna-Vust-theo} F. Knop - “The Luna-Vust theory of spherical embeddings” , Proceedings of the Hyderabad Conference on Algebraic Groups (Hyderabad, 1989) (1991), p. 225-249
\bibitem{Kn1993Uber-Hilberts-vier} F. Knop - “Über Hilberts vierzehntes Problem für Varietäten mit Kompliziertheit eins” , Math. Z. 213 (1993) no. 1, p. 33-36
\bibitem{Kr1984Geometrische-Metho} H. Kraft - Geometrische Methoden in der Invariantentheorie , Aspects of Math. , vol. D1 , Friedr. Vieweg \& Sohn, Braunschweig, 1984
\bibitem{Kraft2017} H. Kraft - “Automorphism groups of affine varieties and a characterization of affine $n$-space” , Trans. Moscow Math. Soc. 78 (2017), p. 171-186
\bibitem{KrReSa2018Is-the-affine-spac} H. Kraft, A. Regeta \& I. van Santen - “Is the affine space determined by its automorphism group?” , Internat. Math. Res. Notices (2019), article ID rny281, 21 pages
\bibitem{Li2010Affine-Bbb-T-varie} A. Liendo - “Affine $\mathbb{T}$-varieties of complexity one and locally nilpotent derivations” , Transform. Groups 15 (2010) no. 2, p. 389-425
\bibitem{Lo2009Proof-of-the-Knop-} I. V. Losev - “Proof of the Knop conjecture” , Ann. Inst. Fourier (Grenoble) 59 (2009) no. 3, p. 1105-1134
\bibitem{Lo2009Uniqueness-propert} I. V. Losev - “Uniqueness property for spherical homogeneous spaces” , Duke Math. J. 147 (2009) no. 2, p. 315-343
\bibitem{LiReUr2018Characterization-o} A. Liendo, A. Regeta \& C. Urech - “Characterization of affine surfaces with a torus action by their automorphism groups” , 2019
\bibitem{Lu2001Varietes-spherique} D. Luna - “Variétés sphériques de type $A$” , Publ. Math. Inst. Hautes Études Sci. (2001) no. 94, p. 161-226
\bibitem{Lu2007La-variete-magnifi} D. Luna - “La variété magnifique modèle” , J. Algebra 313 (2007) no. 1, p. 292-319
\bibitem{LuVu1983Plongements-despac} D. Luna \& T. Vust - “Plongements d’espaces homogènes” , Comment. Math. Helv. 58 (1983) no. 2, p. 186-245
\bibitem{Ra1964A-note-on-automorp} C. P. Ramanujam - “A note on automorphism groups of algebraic varieties” , Math. Ann. 156 (1964), p. 25-33
\bibitem{Re2017Characterization-o} A. Regeta - “Characterization of $n$-dimensional normal affine $\mathrm{SL}_n$-varieties” , 2017
\bibitem{Ro1956Some-basic-theorem} M. Rosenlicht - “Some basic theorems on algebraic groups” , Amer. J. Math. 78 (1956), p. 401-443
\bibitem{1994Algebraic-geometry} - Algebraic geometry. IV (I. R. Shafarevich, ed.) , Encyclopaedia of Math. Sciences , vol. 55 , Springer-Verlag, Berlin, 1994
\bibitem{Ti2011Homogeneous-spaces} D. A. Timashev - Homogeneous spaces and equivariant embeddings , Encyclopaedia of Math. Sciences , vol. 138 , Springer, Heidelberg, 2011
\end{thebibliography}
\end{document}
